% =====================================================================
%  papers/shared/annotated_bibliography.tex
%
%  STUB created by T-005 so the scaffold compiles.  T-001 OWNS THIS FILE
%  and will replace it with the full annotated corpus required by AC-004.
%
%  Contract this stub establishes, which T-001 should preserve:
%    * The file is \input inside a \section, so it starts at \subsection
%      level and must not issue \section.
%    * Each source gets its own \subsection* or description item stating
%      (a) what the source does, (b) why it is credible, (c) what this
%      paper takes from it.
%    * Sources that could not be retrieved are flagged \textbf{Unverified:}.
%    * Every key annotated here must exist in shared/refs.bib and be cited
%      somewhere in the text, or BibTeX will not emit it into the
%      reference list.
% =====================================================================

The reference list that follows the annotations is generated by \textsc{BibTeX}
from \texttt{papers/shared/refs.bib}. Each entry below states what the source
does, why it is credible, and precisely what the present paper takes from it.

\begin{description}\setlength{\itemsep}{0.6em}

\item[\citet{moran1950}] \emph{What it does:} introduces the autocorrelation
  coefficient now universally called Moran's $\MoranI$, as a way of detecting
  whether a quantity measured over a region varies smoothly or erratically in
  space. \emph{Why it is credible:} the founding paper for the statistic, in
  \emph{Biometrika}, and the definition in use today is still Moran's.
  \emph{What we take:} the definition of the statistic and its ratio form.

\item[\citet{getisord1992}] \emph{What it does:} defines the family of distance
  statistics $G_i$ and $\Gistar$ for detecting local concentrations of high or
  low values. \emph{Why it is credible:} the originating paper in
  \emph{Geographical Analysis}; the statistic is implemented under this
  definition in every major spatial toolkit. \emph{What we take:} the definition
  of $\Gistar$, including the fact that it standardises against the
  \emph{global} mean and standard deviation.

\item[\citet{anselin1995}] \emph{What it does:} sets out the LISA framework,
  which decomposes a global measure of spatial association into per-location
  contributions. \emph{Why it is credible:} the standard reference for local
  spatial statistics, and the basis of the local Moran implementations in
  \texttt{spdep}, \texttt{esda} and ArcGIS. \emph{What we take:} the framing of
  local statistics as decompositions of a global one.

\end{description}

\begin{keyinsight}[Placeholder]
This is a stub annotated bibliography supplied by the build scaffold (T-005) to
prove the \textsc{BibTeX} wiring. T-001 replaces this file, and
\texttt{refs.bib}, with the retrieved and verified corpus required by AC-004 and
AC-005.
\end{keyinsight}
